% Fixture: an amsart paper with equation numbers on the left, alphabetic citation labels and a two-column
% reference list. Compile with: tectonic leqno-alpha.tex
\documentclass[leqno]{amsart}
\usepackage{amssymb}
\newtheorem{theorem}{Theorem}[section]
\newtheorem{lemma}[theorem]{Lemma}
\begin{document}
\title{Recovery of Sparse Signals}
\author{A. Fixture}
\maketitle
\section{Setting}
Let $A$ be an $n \times m$ measurement matrix with $n < m$, and consider the program
\begin{equation}\label{eq:p1}
\min_x \|x\|_{\ell_1} \quad \text{subject to} \quad Ax = y .
\end{equation}
Program \eqref{eq:p1} is known as basis pursuit [CDS98]. Its stability was studied in [CRT06, Theorem 1].
\begin{equation}\label{eq:rip}
(1-\delta_S)\|c\|^2 \le \|A_T c\|^2 \le (1+\delta_S)\|c\|^2 .
\end{equation}
The inequality \eqref{eq:rip} is the restricted isometry property [CT05].
\section{Main result}
\begin{theorem}\label{thm:main}
If $\delta_{3S} + 3\delta_{4S} < 2$, then the solution of \eqref{eq:p1} recovers every $S$-sparse signal.
\end{theorem}
\begin{proof}
Combine \eqref{eq:rip} with Lemma \ref{lem:cone} below.
\end{proof}
\begin{lemma}\label{lem:cone}
The error $h = x^\sharp - x_0$ obeys the cone constraint $\|h_{T^c}\|_{\ell_1} \le \|h_T\|_{\ell_1}$.
\end{lemma}
\begin{proof}
This follows from the optimality of $x^\sharp$ in \eqref{eq:p1}; see also [Don06].
\end{proof}
\begin{thebibliography}{CRT06}
\bibitem[CDS98]{cds} S. S. Chen, D. L. Donoho and M. A. Saunders, Atomic decomposition by basis pursuit, SIAM J. Sci. Comput. 20 (1998), 33--61.
\bibitem[CRT06]{crt} E. J. Cand\`es, J. Romberg and T. Tao, Stable signal recovery from incomplete and inaccurate measurements, Comm. Pure Appl. Math. 59 (2006), 1207--1223.
\bibitem[CT05]{ct} E. J. Cand\`es and T. Tao, Decoding by linear programming, IEEE Trans. Inform. Theory 51 (2005), 4203--4215.
\bibitem[Don06]{don} D. L. Donoho, Compressed sensing, IEEE Trans. Inform. Theory 52 (2006), 1289--1306.
\end{thebibliography}
\end{document}
